%% kapa: exam=1 round=31 year=2020 period=1 form=A subject=economics subject_ko=경제학원론 seq=40 number=80
%% source: Q-Net 감정평가사 기출문제 (cst00309) artlSeq=5207219 file=extracted/1008/5207219/A형/A형/2020년도 제31회 감정평가사 1차 시험_1교시_A형.pdf
%% conversion: kordoc 4.24.51 -> tools/gen1.py (auto); review=proofread-sonnet; flags=none
\begin{question}{80}
리카디언 등가(Ricardian equivalence) 정리에 관한 설명으로 옳지 않은 것은?
\begin{choices}
  \item 민간 경제주체는 합리적 기대를 한다.
  \item 소비자가 차입 제약에 직면하면 이 정리는 성립되지 않는다.
  \item 소비자가 근시안적 견해를 가지면 이 정리는 성립되지 않는다.
  \item 현재의 감세가 현재의 민간소비를 증가시킨다는 주장과는 상반된 것이다.
  \item 정부가 미래의 정부지출을 축소한다는 조건에서 현재 조세를 줄이는 경우에 현재의 민간소비는 변하지 않는다.
\end{choices}
\end{question}
\answer{5}
